% Generated draft: requires mathematical and source review.
\documentclass[11pt]{article}
\usepackage[margin=25mm]{geometry}
\usepackage{fontspec}
\setmainfont{Latin Modern Roman}
\usepackage{amsmath,amssymb,mathtools,unicode-math}
\setmathfont{Latin Modern Math}
\usepackage{enumitem,needspace,xurl,hyperref}
\setlength{\parindent}{0pt}
\setlength{\parskip}{5pt}
\setlength{\emergencystretch}{3em}
\newcommand{\N}{\mathbb{N}}
\newcommand{\Z}{\mathbb{Z}}
\newcommand{\Q}{\mathbb{Q}}
\newcommand{\R}{\mathbb{R}}
\newcommand{\C}{\mathbb{C}}
\newcommand{\F}{\mathbb{F}}
\newcommand{\A}{\mathbb{A}}
\newcommand{\PP}{\mathbb{P}}
\newcommand{\E}{\mathbb{E}}
\newcommand{\eps}{\varepsilon}
\newcommand{\dd}{\,\mathrm{d}}
\begin{document}
\section*{UnsolvedMath ID: 700853}
\textbf{Source problem code:} catalogue entry
\subsection{Definitions and mathematical statement}
Let \(G\) be a finite group, let \(|G|\) denote its order, and let \(\operatorname{Irr}_{\C}(G)\) denote the set of all irreducible complex characters of \(G\), including the trivial character. For \(\chi\in\operatorname{Irr}_{\C}(G)\), its degree \(\chi(1)\) is the dimension of the corresponding irreducible complex representation. Each irreducible character is counted once. Describe, up to group isomorphism, all finite groups \(G\) satisfying
\[
\sum_{\chi\in\operatorname{Irr}_{\C}(G)}\chi(1)^3
\leq |G|\log_2|G|.
\]
Here \(\log_2\) is the logarithm to base \(2\). No restriction on the order of \(G\) beyond finiteness is stated.
\subsection{Short English statement}
Describe the finite groups for which the sum of the cubes of all irreducible complex character degrees is at most the group order times its base-\(2\) logarithm. The source identifies applications in signal processing as motivation.
\subsection{Sources}
\begin{itemize}
\item Local UnsolvedMath registry, permanent ID \(700853\), exact source code \texttt{catalogue entry}, titled \emph{Kourovka Notebook Problem 12.34}; supplied upstream revision \texttt{d371ee917ac9647fbb06cd6fc529deec6f3dda82}.
\item L. S. Kazarin, Problem 12.34, \emph{The Kourovka Notebook (Unsolved Problems in Group Theory)}, 12th issue (1992). The supplied record also identifies the 21st issue edition (2026), \texttt{arXiv:1401.0300v48}. The collection description names E. I. Khukhro and V. D. Mazurov as editors.
\item The supplied registry labels the problem \emph{open}. This status and the bibliographic identification are supplied source information, not independently verified assessments.
\end{itemize}
\end{document}
